
\documentclass{article}
\usepackage{colortbl}
\usepackage{makecell}
\usepackage{multirow}
\usepackage{supertabular}

\begin{document}

\newcounter{utterance}

\centering \large Interaction Transcript for game `wordle', experiment `medium\_frequency\_words\_no\_clue\_no\_critic', episode 1 with mistral{-}large{-}3{-}azure.
\vspace{24pt}

{ \footnotesize  \setcounter{utterance}{1}
\setlength{\tabcolsep}{0pt}
\begin{supertabular}{c@{$\;$}|p{.15\linewidth}@{}p{.15\linewidth}p{.15\linewidth}p{.15\linewidth}p{.15\linewidth}p{.15\linewidth}}
    \# & \multicolumn{2}{c}{Player} && \multicolumn{2}{c}{Game Master} \\
    \hline

    \theutterance \stepcounter{utterance}  
    & & & \multicolumn{4}{p{0.6\linewidth}}{
        \cellcolor[rgb]{0.9,0.9,0.9}{
            \makecell[{{p{\linewidth}}}]{
                \texttt{\tiny{[P1$\langle$GM]}}
                \texttt{Vi ste jezički čarobnjak koji voli da pogađa reči koristeći data pravila.} \\
\\ 
\texttt{Dobrodošli u Wordle! Imate šest pokušaja da pogodite ciljnu reč, važeću reč od pet malih slova, napisanu latinicom. Molimo koristite oznake "objašnjenje:" i "pogađanje:" da biste dali sažeto objašnjenje za svako pogađanje.} \\
\\ 
\texttt{Na primer, ako je vaše pogađanje "vreme", vaš odgovor treba da bude} \\
\texttt{objašnjenje: ovo je česta reč od pet slova i započinjem pogađanje ovom rečju.} \\
\texttt{pogađanje: vreme} \\
\\ 
\texttt{Nakon svakog pogađanja, vaš odgovor će biti proveren i dobićete povratnu informaciju koja pokazuje koja su slova tačna (zeleno), koja su slova tačna ali na pogrešnom mestu (žuto) i koja su slova netačna (crveno). Ova povratna informacija može biti korisna pri određivanju koja slova da uključite ili isključite u sledećem pogađanju.} \\
\\ 
\texttt{Na primer, povratna informacija za "vreme" može biti:} \\
\texttt{povratna\_informacija\_o\_pogađanju: v<žuto> r<žuto> e<zeleno> m<žuto> e<crveno>} \\
\\ 
\texttt{Objašnjenje treba da sadrži detalje o tome kako se povratna\_informacija\_o\_pogađanju koristi da bi se došlo do novog pogađanja.} \\
\\ 
\texttt{Ne generišite nikakav drugi tekst.} \\
\\ 
\texttt{Počnimo sa vašim prvim pogađanjem.} \\
            }
        }
    }
     \\ \\

    \theutterance \stepcounter{utterance}  
    & \multicolumn{4}{p{0.6\linewidth}}{
        \cellcolor[rgb]{0.561,0.737,0.561}{
            \makecell[{{p{\linewidth}}}]{
                \texttt{\tiny{[P1$\rangle$GM]}}
                \texttt{objašnjenje: Počinjem sa rečju "slova" jer sadrži često korišćena slova (s, l, o, v, a) koja će mi pomoći da brzo eliminišem ili potvrdim njihovu prisutnost u ciljnoj reči.} \\
\texttt{pogađanje: slova} \\
            }
        }
    }
    & & \\ \\

    \theutterance \stepcounter{utterance}  
    & & & \multicolumn{2}{p{0.3\linewidth}}{
        \cellcolor[rgb]{0.9,0.9,0.9}{
            \makecell[{{p{\linewidth}}}]{
                \texttt{\tiny{[GM$|$GM]}}
                \texttt{attempts = 1} \\
\texttt{target\_word = rizik} \\
\texttt{guess = slova} \\
\texttt{guess\_feedback = s<red> l<red> o<red> v<red> a<red>} \\
            }
        }
    }
    & & \\ \\

    \theutterance \stepcounter{utterance}  
    & & & \multicolumn{4}{p{0.6\linewidth}}{
        \cellcolor[rgb]{0.9,0.9,0.9}{
            \makecell[{{p{\linewidth}}}]{
                \texttt{\tiny{[P1$\langle$GM]}}
                \texttt{povratna\_informacija\_o\_pogađanju: s<crveno> l<crveno> o<crveno> v<crveno> a<crveno>} \\
\\ 
\texttt{Navedite svoj odgovor samo u ovom formatu.} \\
\texttt{objašnjenje: detalji} \\
\texttt{pogađanje: reč} \\
            }
        }
    }
     \\ \\

    \theutterance \stepcounter{utterance}  
    & \multicolumn{4}{p{0.6\linewidth}}{
        \cellcolor[rgb]{0.561,0.737,0.561}{
            \makecell[{{p{\linewidth}}}]{
                \texttt{\tiny{[P1$\rangle$GM]}}
                \texttt{objašnjenje: Sva slova iz prvog pokušaja ("slova") su netačna (crvena), pa ih isključujem. Sada biram reč "truni" jer sadrži nova, često korišćena slova (t, r, u, n, i) koja nisu testirana u prethodnom pokušaju.} \\
\texttt{pogađanje: truni} \\
            }
        }
    }
    & & \\ \\

    \theutterance \stepcounter{utterance}  
    & & & \multicolumn{2}{p{0.3\linewidth}}{
        \cellcolor[rgb]{0.9,0.9,0.9}{
            \makecell[{{p{\linewidth}}}]{
                \texttt{\tiny{[GM$|$GM]}}
                \texttt{Error: The guessed word is not a valid word for this game.} \\
            }
        }
    }
    & & \\ \\

    \theutterance \stepcounter{utterance}  
    & & & \multicolumn{4}{p{0.6\linewidth}}{
        \cellcolor[rgb]{0.9,0.9,0.9}{
            \makecell[{{p{\linewidth}}}]{
                \texttt{\tiny{[P1$\langle$GM]}}
                \texttt{Vaše pogađanje nije važeća reč za ovu igru. Molimo pokušajte ponovo.} \\
\\ 
\texttt{Navedite svoj odgovor samo u ovom formatu.} \\
\texttt{objašnjenje: detalji} \\
\texttt{pogađanje: reč} \\
            }
        }
    }
     \\ \\

    \theutterance \stepcounter{utterance}  
    & \multicolumn{4}{p{0.6\linewidth}}{
        \cellcolor[rgb]{0.561,0.737,0.561}{
            \makecell[{{p{\linewidth}}}]{
                \texttt{\tiny{[P1$\rangle$GM]}}
                \texttt{objašnjenje: Pošto "truni" nije validna reč, biram drugu čestu reč sa novim slovima: "kamen". Ova reč sadrži slova koja nisu korišćena u prethodnom pokušaju (k, a, m, e, n) i omogućava mi da testiram njihovu prisutnost.} \\
\texttt{pogađanje: kamen} \\
            }
        }
    }
    & & \\ \\

    \theutterance \stepcounter{utterance}  
    & & & \multicolumn{2}{p{0.3\linewidth}}{
        \cellcolor[rgb]{0.9,0.9,0.9}{
            \makecell[{{p{\linewidth}}}]{
                \texttt{\tiny{[GM$|$GM]}}
                \texttt{attempts = 2} \\
\texttt{target\_word = rizik} \\
\texttt{guess = kamen} \\
\texttt{guess\_feedback = k<yellow> a<red> m<red> e<red> n<red>} \\
            }
        }
    }
    & & \\ \\

    \theutterance \stepcounter{utterance}  
    & & & \multicolumn{4}{p{0.6\linewidth}}{
        \cellcolor[rgb]{0.9,0.9,0.9}{
            \makecell[{{p{\linewidth}}}]{
                \texttt{\tiny{[P1$\langle$GM]}}
                \texttt{povratna\_informacija\_o\_pogađanju: k<žuto> a<crveno> m<crveno> e<crveno> n<crveno>} \\
\\ 
\texttt{Navedite svoj odgovor samo u ovom formatu.} \\
\texttt{objašnjenje: detalji} \\
\texttt{pogađanje: reč} \\
            }
        }
    }
     \\ \\

    \theutterance \stepcounter{utterance}  
    & \multicolumn{4}{p{0.6\linewidth}}{
        \cellcolor[rgb]{0.561,0.737,0.561}{
            \makecell[{{p{\linewidth}}}]{
                \texttt{\tiny{[P1$\rangle$GM]}}
                \texttt{objašnjenje: Iz povratne informacije vidim da je samo slovo **k** prisutno u reči, ali na pogrešnom mestu (žuto), dok su **a, m, e, n** netačna (crvena). Sledeća reč bi trebalo da zadrži **k** (ali na drugoj poziciji) i uvede nova slova koja nisu testirana. Biram "kisel" jer sadrži **k** na drugom mestu i nova slova (i, s, e, l).} \\
\texttt{pogađanje: kisel} \\
            }
        }
    }
    & & \\ \\

    \theutterance \stepcounter{utterance}  
    & & & \multicolumn{2}{p{0.3\linewidth}}{
        \cellcolor[rgb]{0.9,0.9,0.9}{
            \makecell[{{p{\linewidth}}}]{
                \texttt{\tiny{[GM$|$GM]}}
                \texttt{attempts = 3} \\
\texttt{target\_word = rizik} \\
\texttt{guess = kisel} \\
\texttt{guess\_feedback = k<yellow> i<green> s<red> e<red> l<red>} \\
            }
        }
    }
    & & \\ \\

    \theutterance \stepcounter{utterance}  
    & & & \multicolumn{4}{p{0.6\linewidth}}{
        \cellcolor[rgb]{0.9,0.9,0.9}{
            \makecell[{{p{\linewidth}}}]{
                \texttt{\tiny{[P1$\langle$GM]}}
                \texttt{povratna\_informacija\_o\_pogađanju: k<žuto> i<zeleno> s<crveno> e<crveno> l<crveno>} \\
\\ 
\texttt{Navedite svoj odgovor samo u ovom formatu.} \\
\texttt{objašnjenje: detalji} \\
\texttt{pogađanje: reč} \\
            }
        }
    }
     \\ \\

    \theutterance \stepcounter{utterance}  
    & \multicolumn{4}{p{0.6\linewidth}}{
        \cellcolor[rgb]{0.561,0.737,0.561}{
            \makecell[{{p{\linewidth}}}]{
                \texttt{\tiny{[P1$\rangle$GM]}}
                \texttt{objašnjenje: Iz povratne informacije:} \\
\texttt{{-} **i** je tačno na svom mestu (zeleno, 2. slovo).} \\
\texttt{{-} **k** je prisutno, ali na pogrešnom mestu (žuto), pa mora biti negde osim na 1. poziciji.} \\
\texttt{{-} **s, e, l** su netačna (crvena) i isključujem ih.} \\
\\ 
\texttt{Sledeća reč treba da:} \\
\texttt{1. Ima **i** na 2. mestu.} \\
\texttt{2. Sadrži **k**, ali ne na 1. mestu.} \\
\texttt{3. Koristi nova slova (npr. **d, o, r, t, p**).} \\
\\ 
\texttt{Biram "krilo" jer:} \\
\texttt{{-} **i** je na 2. mestu.} \\
\texttt{{-} **k** je na 1. mestu (što je pogrešno, ali testiram), a **r, l, o** su nova slova.} \\
\texttt{*(Napomena: Ako "krilo" ne bude valjano, sledeća opcija bi bila "pirka" ili "tirka".)*} \\
\\ 
\texttt{pogađanje: krilo} \\
            }
        }
    }
    & & \\ \\

    \theutterance \stepcounter{utterance}  
    & & & \multicolumn{2}{p{0.3\linewidth}}{
        \cellcolor[rgb]{0.9,0.9,0.9}{
            \makecell[{{p{\linewidth}}}]{
                \texttt{\tiny{[GM$|$GM]}}
                \texttt{attempts = 4} \\
\texttt{target\_word = rizik} \\
\texttt{guess = krilo} \\
\texttt{guess\_feedback = k<yellow> r<yellow> i<yellow> l<red> o<red>} \\
            }
        }
    }
    & & \\ \\

    \theutterance \stepcounter{utterance}  
    & & & \multicolumn{4}{p{0.6\linewidth}}{
        \cellcolor[rgb]{0.9,0.9,0.9}{
            \makecell[{{p{\linewidth}}}]{
                \texttt{\tiny{[P1$\langle$GM]}}
                \texttt{povratna\_informacija\_o\_pogađanju: k<žuto> r<žuto> i<žuto> l<crveno> o<crveno>} \\
\\ 
\texttt{Navedite svoj odgovor samo u ovom formatu.} \\
\texttt{objašnjenje: detalji} \\
\texttt{pogađanje: reč} \\
            }
        }
    }
     \\ \\

    \theutterance \stepcounter{utterance}  
    & \multicolumn{4}{p{0.6\linewidth}}{
        \cellcolor[rgb]{0.561,0.737,0.561}{
            \makecell[{{p{\linewidth}}}]{
                \texttt{\tiny{[P1$\rangle$GM]}}
                \texttt{objašnjenje:} \\
\texttt{{-} **i** je ponovo žuto (a ranije je bilo zeleno na 2. mestu), što znači da **i** **nije** na 2. poziciji – greška u prethodnoj analizi! Ciljna reč **ne sadrži i na 2. mestu**, ali sadrži **i** negde drugo.} \\
\texttt{{-} **k** i **r** su žuti (prisutni, ali na pogrešnom mestu).} \\
\texttt{{-} **l** i **o** su crvena (isključena).} \\
\\ 
\texttt{Nova reč treba da:} \\
\texttt{1. **Ne sadrži i na 2. mestu** (ranije je bilo zeleno greškom, sada je žuto).} \\
\texttt{2. Sadrži **k** i **r**, ali ne na pozicijama iz prethodnih pokušaja (k≠1, r≠2).} \\
\texttt{3. Izbegava **s, e, l, o, a, m, n** (crvena).} \\
\texttt{4. Testira nova slova (npr. **d, t, p, u, j**).} \\
\\ 
\texttt{Biram **"krump"** (ako nije validno, sledeća opcija: **"krupi"**):} \\
\texttt{{-} **k** na 1. mestu (testiram da li je žuto zbog pozicije ili slova).} \\
\texttt{{-} **r** na 2. mestu (isto test).} \\
\texttt{{-} **u** i **p** su nova slova.} \\
\\ 
\texttt{pogađanje: krupi} \\
            }
        }
    }
    & & \\ \\

    \theutterance \stepcounter{utterance}  
    & & & \multicolumn{2}{p{0.3\linewidth}}{
        \cellcolor[rgb]{0.9,0.9,0.9}{
            \makecell[{{p{\linewidth}}}]{
                \texttt{\tiny{[GM$|$GM]}}
                \texttt{attempts = 5} \\
\texttt{target\_word = rizik} \\
\texttt{guess = krupi} \\
\texttt{guess\_feedback = k<yellow> r<yellow> u<red> p<red> i<yellow>} \\
            }
        }
    }
    & & \\ \\

    \theutterance \stepcounter{utterance}  
    & & & \multicolumn{4}{p{0.6\linewidth}}{
        \cellcolor[rgb]{0.9,0.9,0.9}{
            \makecell[{{p{\linewidth}}}]{
                \texttt{\tiny{[P1$\langle$GM]}}
                \texttt{povratna\_informacija\_o\_pogađanju: k<žuto> r<žuto> u<crveno> p<crveno> i<žuto>} \\
\\ 
\texttt{Navedite svoj odgovor samo u ovom formatu.} \\
\texttt{objašnjenje: detalji} \\
\texttt{pogađanje: reč} \\
            }
        }
    }
     \\ \\

    \theutterance \stepcounter{utterance}  
    & \multicolumn{4}{p{0.6\linewidth}}{
        \cellcolor[rgb]{0.561,0.737,0.561}{
            \makecell[{{p{\linewidth}}}]{
                \texttt{\tiny{[P1$\rangle$GM]}}
                \texttt{objašnjenje:} \\
\texttt{{-} **k, r, i** su žuti (prisutni, ali na pogrešnim mestima):} \\
\texttt{  {-} **k** ≠ 1. pozicija (iz "krupi").} \\
\texttt{  {-} **r** ≠ 2. pozicija (iz "krupi").} \\
\texttt{  {-} **i** ≠ 5. pozicija (iz "krupi").} \\
\texttt{{-} **u, p** su crvena (isključena).} \\
\texttt{{-} Ranije smo videli da **i** nije na 2. mestu (iz "kisel" → žuto).} \\
\\ 
\texttt{Nove informacije:} \\
\texttt{1. Reč **mora sadržati k, r, i** (ali na drugim pozicijama).} \\
\texttt{2. Isključena slova: **s, l, o, a, m, n, e, v, u, p**.} \\
\texttt{3. Moguća pozicija za **i**: 1, 3, 4 (ne 2, ne 5).} \\
\texttt{4. Moguća pozicija za **k** i **r**: 3, 4, 5 (ne 1, ne 2).} \\
\\ 
\texttt{Biram **"rikša"** (ako nije validno, sledeća opcija: **"kirka"**):} \\
\texttt{{-} **r** na 1. mestu (testiram novu poziciju).} \\
\texttt{{-} **i** na 2. mestu (testiram, iako je ranije bilo žuto na 2. mestu – možda greška u analizi).} \\
\texttt{{-} **k** na 3. mestu (testiram).} \\
\texttt{{-} **š** i **a** su nova slova (a je već isključeno, ali testiram š).} \\
\\ 
\texttt{pogađanje: kirka} \\
            }
        }
    }
    & & \\ \\

    \theutterance \stepcounter{utterance}  
    & & & \multicolumn{2}{p{0.3\linewidth}}{
        \cellcolor[rgb]{0.9,0.9,0.9}{
            \makecell[{{p{\linewidth}}}]{
                \texttt{\tiny{[GM$|$GM]}}
                \texttt{attempts = 6} \\
\texttt{target\_word = rizik} \\
\texttt{guess = kirka} \\
\texttt{guess\_feedback = k<yellow> i<green> r<yellow> k<red> a<red>} \\
            }
        }
    }
    & & \\ \\

    \theutterance \stepcounter{utterance}  
    & & & \multicolumn{2}{p{0.3\linewidth}}{
        \cellcolor[rgb]{0.9,0.9,0.9}{
            \makecell[{{p{\linewidth}}}]{
                \texttt{\tiny{[GM$|$GM]}}
                \texttt{game\_result = LOSS} \\
            }
        }
    }
    & & \\ \\

\end{supertabular}
}

\end{document}
